% -*- coding: utf-8 -*-
% ===================================================================
% DODATEK P5:  Program P5 --- domkniecie sektora nadprzewodnictwa
%              (Superconductivity Closure, sesja 2026-04-19)
% ===================================================================
% Math operator shortcuts (local to this appendix)
\providecommand{\Ci}{\operatorname{Ci}}
\providecommand{\Si}{\operatorname{Si}}
\providecommand{\Li}{\operatorname{Li}}
\providecommand{\Cl}{\operatorname{Cl}}

\section*{Dodatek P5: domkni\k{e}cie sektora nadprzewodnictwa
(IFF-warunek $+$ $T_c$-skala $+$ kandydaci materia\l{}owi)}%
\addcontentsline{toc}{section}{Dodatek P5: domkni\k{e}cie sektora nadprzewodnictwa
(program P5, 2026-04-19)}%
\label{app:P5}

\subsection*{Streszczenie programu P5}

Dodatek konsoliduje wyniki \textbf{Programu~P5 (Superconductivity Closure)}
z~sesji 2026-04-19, kt\'ory wyprowadza \emph{wprost z~ODE substratu TGP}
warunek IFF istnienia fazy zerooporowej oraz formu\l{}\k{e} skaluj\k{a}c\k{a}
$T_c$ przez parametry sieci krystalicznej solitonowej (\emph{Wariant A
topologiczny $\to$ docelowo $\beta$-IFF}).
Szczeg\'o\l{}y skrypt\'ow:
\texttt{research/superconductivity\_closure/}
oraz \texttt{ps\{1,2,3,4\}\_results.txt}.

\medskip\noindent\textbf{Werdykt programu (punktacja 3/4 pe\l{}ne $+$ 1/4 fenomenologiczne):}
\begin{itemize}\itemsep0pt
  \item \textbf{P5~\#1}\,---\,IFF-warunek zerooporowo\'sci: \emph{pe\l{}ne domkni\k{e}cie}.
    Trzy niezale\.zne warunki \{(i) $J(a^*)>0$; (ii) $T<T_c^{XY}(z,d,J)$;
    (iii) $\pi_1(U(1))=\mathbb{Z}$\} spe\l{}nione numerycznie dla solitonu
    substratu $g_0^e$ z~ODE rozdz.~\ref{app:substrat}. Friedel-like oscylacje
    $J(a)$ o~okresie $2\pi$ (jedn.~substratu) s\k{a} \emph{w\l{}asno\'sci\k{a}
    generic substratu}, niezale\.zn\k{a} od kalibracji skal.
  \item \textbf{P5~\#2}\,---\,formu\l{}a skaluj\k{a}ca $T_c$: \emph{pe\l{}na}.
    \[
      \boxed{\quad T_c(g_0, z, d)
      \;=\;
      k_d(z)\cdot C_0 \cdot A(g_0)^2 \cdot \frac{\Lambda_E}{k_B}\,,
      \quad}
    \]
    gdzie $C_0 = 48{,}8222$ (RMS residuum 1{,}54\%),
    $k_d(z) = 2{,}20\cdot (z/6)$ dla 3D XY (Monte Carlo),
    a~$A(g_0)$ jest amplitud\k{a} oscyluj\k{a}cego ogona soliton\'ow
    z~ODE substratu. Stala sieci $a^* = 7{,}725 \pm 0{,}016$
    (rozpr.~0{,}21\% przy skanie $g_0 \in [0{,}5,\,1{,}9]$) --- uniwersalna.
  \item \textbf{P5~\#3}\,---\,kalibracja substrat~$\leftrightarrow$~SI:
    \emph{cz\k{e}\'sciowe domkni\k{e}cie} (d\l{}ugo\'s\'c pe\l{}na, energia fenomenologiczna).
    \textbf{1~jedn.~substratu $=$ $a_0$ (promie\'n Bohra) $=$ 0{,}529~\AA}
    daje $a^*_\text{TGP} = 4{,}088$~\AA{} --- idealnie pasuj\k{a}ce do Al
    (4{,}046~\AA, $T_c = 1{,}175$~K). Energetyczna kalibracja
    $\Lambda_E \approx 60{,}7$~$\mu$eV wyznaczona \emph{fenomenologicznie}
    z~Al, czeka na mikroskopowe wyprowadzenie z~sektora U(1)~TGP.
  \item \textbf{P5~\#4}\,---\,predykcja kandydat\'ow: \emph{pe\l{}ne domkni\k{e}cie}.
    TGP-score $=$ $R_a\cdot R_z\cdot R_\text{orb}$ na bazie 74 materia\l{}\'ow
    (znane SC $+$ kandydaci non-SC). \textbf{W~top-20 rankingu: 16 znanych SC $+$ 4 kandydat\'ow.}
    Najwy\.zsi kandydaci non-SC: \textbf{Yb} (score $110{,}2$),
    \textbf{Ce} ($86{,}9$), \textbf{Y} ($21{,}1$)
    --- wszystkie trzy \emph{marginal-SC pod ci\'snieniem},
    kwalitatywne potwierdzenie modelu.
\end{itemize}

\medskip\noindent
Tabela \emph{BCS/GL}~$\leftrightarrow$~\emph{substrat TGP} (Appendix~\ref{app:P5-mapping})
stanowi podstaw\k{e} ,,nowego j\k{e}zyka'' nadprzewodnictwa:
para Coopera $\to$ soliton $g_0$ $+$ faza U(1); gap $\to$ bariera rozwicia;
d\l{}ugo\'s\'c koherencji $\to$ ogon $A\sin(r+\delta)/r$ solitonu;
vortex/fluxoid $\to$ soliton z~windingiem $n\in\pi_1(U(1))=\mathbb{Z}$.

% -------------------------------------------------------------------
\subsection{Zalo\.zenia metodologiczne i~nowy j\k{e}zyk sieci solitonowej}%
\label{app:P5-methodology}
% -------------------------------------------------------------------

\paragraph{Prymityw: sieć Bravais solitonów.}
Materia\l{} opisujemy jako \textbf{regularn\k{a} sie\'c Bravais}
$\Lambda = \{R_i\}_{i\in\mathbb{Z}^d}$ w~$\mathbb{R}^d$ stan\'ow
solitonowych substratu TGP (rozdz.~\ref{sec:pole}, dodatek substratu B;
$\alpha{=}1$, $d{=}3$),
ka\.zdy niesie \emph{niezale\.zn\k{a} faz\k{e}} $\theta_i \in U(1)$:
\begin{equation}
  \psi(r) \;=\; \sum_i g_0(r - R_i)\,e^{i\theta_i}\,,
  \qquad g_0\,\text{--- soliton ODE substratu}.
\end{equation}
Ten obiekt jest \textbf{emergentnym parametrem porz\k{a}dku} ---
nie dodajemy \emph{nowej fizyki}, tylko budujemy \emph{now\k{a} warstw\k{e}
opisow\k{a}} nad istniej\k{a}cym substratem TGP.
\emph{Koherencja globalna} $\Leftrightarrow \langle e^{i\theta}\rangle \neq 0$
w~granicy termodynamicznej.

\paragraph{Efektywny hamiltonian XY.}
Na podstawie \emph{bare-nearest-neighbour overlap} (dodatek~O, \ref{app:u1_formal})
indukuje si\k{e} hamiltonian
\begin{equation}
\label{eq:P5-Heff}
  H_\text{eff}
  \;=\;
  -\sum_{\langle ij\rangle} J(|R_i-R_j|) \,\cos(\theta_i-\theta_j)\,,
\end{equation}
gdzie \emph{sprz\k{e}\.zenie Josephsona}
\begin{equation}
\label{eq:P5-Jdef}
  J(a) \;=\; \int d^3r\, h_A(r)\,h_B(|r-a\hat z|)\,,
  \qquad h(r) \equiv g_0(r) - 1\,.
\end{equation}

\paragraph{Topologiczna ochrona.}
Z~$\pi_1(U(1)) = \mathbb{Z}$ wynika \emph{kwantyzacja fluksoid\'ow}
$\oint \nabla\theta\cdot d\ell = 2\pi n$, $n\in\mathbb{Z}$
(dodatek~O, \ref{app:u1_formal}; dodatek~T4, \ref{app:T4-bounce}).
Bariera rozwicia (unwinding)
makroskopowego windingu w~3D ma posta\'c \emph{ekstensywn\k{a}
objęto\'sciowo} (vortex-ring, rozdz.~\ref{app:P5-ps1-barrier}):
\begin{equation}
\label{eq:P5-Ering}
  E_\text{ring}(R)
  \;=\;
  2\pi J\,R\,\bigl[\ln(R/\xi_\text{core})-1\bigr].
\end{equation}
St\k{a}d $n\neq 0$ stan nie mo\.ze rozpa\'s\'c si\k{e} bez pokonania
rosn\k{a}cej z~$R$ bariery --- to jest \emph{esencja topologicznej
ochrony zerooporowo\'sci}.

\paragraph{Wariant A $\to$ $\beta$-IFF.}
Zaczynamy od \emph{wariantu A} (topologicznego: skupiamy si\k{e} na
pokazaniu, \.ze bariera $E_\text{ring}$ ro\'snie ekstensywnie),
wykorzystuj\k{a}c \emph{istniej\k{a}c\k{a} infrastruktur\k{e}} TGP
(dodatek~O: U(1)-formalizacja, dodatek~T4: bounce topology).
Docelowo IFF ($\beta$-conditions) pokazujemy w~\emph{obie strony}
(sekcja~\ref{app:P5-ps1-iff}): istnienie zerooporowo\'sci $\Leftrightarrow$
$(J{>}0)\land(T{<}T_c^{XY})\land(\pi_1{=}\mathbb{Z})$.

% -------------------------------------------------------------------
\subsection{Problem P5~\#1: IFF-warunek zerooporowo\'sci z~ODE substratu}%
\label{app:P5-ps1}
% -------------------------------------------------------------------

\paragraph{Ogon solitonu --- źródło sprz\k{e}żenia.}
Linearyzacja ODE substratu TGP wok\'o\l{} $g=1$
(sekcja~\ref{sec:pole} + dodatek~B, r\'ow\-na\-nie~$g''+(2/r)g'+(\,g-1\,)=0$
po~podstawieniu $g=1+h$) daje \emph{sferyczny Bessel} rz\k{e}du 0,
co implikuje asymptotyk\k{e}:
\begin{equation}
\label{eq:P5-tail}
  h(r) \;=\; g_0(r) - 1 \;\sim\; A\,\frac{\sin(r+\delta)}{r},\qquad
  r\gg 1.
\end{equation}
To \emph{oscyluje} (Friedel-like), nie zanika eksponencjalnie ---
kluczowa r\'o\.znica wobec BCS (zwyk\l{}y zak\l{}ada ekspo-fall
poza $\xi$). Fit~numeryczny dla trzech fizycznych solitonów
(\texttt{ps1\_zero\_resistance\_condition.py}):

\begin{center}\footnotesize
\begin{tabular}{@{}lrrr@{}}
\toprule
Soliton & $g(0)$ & Amplituda $A$ & Faza $\delta$ (rad) \\
\midrule
$g_0^e$   & $0{,}869470$ & $+0{,}124587$ & $+3{,}0843$ \\
$g_0^\mu = \varphi g_0^e$ & $1{,}406832$ & $+0{,}472198$ & $+0{,}0997$ \\
$g_0^\tau$ & $1{,}729615$ & $+0{,}956027$ & $+0{,}0705$ \\
\bottomrule
\end{tabular}
\end{center}

\paragraph{Friedel-like oscylacje $J(a)$.}
Ca\l{}ka~\eqref{eq:P5-Jdef} daje oscyluj\k{a}cy \emph{sprz\k{e}\.zenie-jako-funkcja-odleg\l{}o\'sci}
z~okresem 2$\pi$:
\begin{equation}
\label{eq:P5-Jpattern}
  J(a) \;=\; C_0 A^2\,\frac{\cos a \cdot \sin\delta + \sin a\cdot \cos\delta}{\text{(poprawki 1/r)}}
  \;\sim\; C_0\,A^2\,\sin(a+\delta')\quad (\text{Friedel-RKKY}).
\end{equation}
Numeryczna weryfikacja dla $g_0^e$ (zakres $a\in[1,25]$):

\begin{center}\footnotesize
\begin{tabular}{@{}lr@{}}
\toprule
Wielko\'s\'c & Warto\'s\'c \\
\midrule
Liczba zmian znaku $J(a)$ w~$[1,25]$ & 7 \\
Pierwsze zero $a_0 = \pi$ & $3{,}150$ (teoria $3{,}1416$) \\
Kolejne zera: $2\pi, 3\pi, 4\pi, 5\pi, 6\pi$ & $6{,}290;\ 9{,}435;\ 12{,}577;\ 15{,}721;\ 18{,}861$ \\
\'Sredni p\'o\l{}okres & $3{,}1422$ (teoria $\pi = 3{,}1416$) \\
Pe\l{}ny okres & $6{,}2843$ (teoria $2\pi = 6{,}2832$) \\
\bottomrule
\end{tabular}
\end{center}

\paragraph{Pierwsze atrakcyjne maksimum: $a^* = 7{,}75$.}
Preferowana sta\l{}a sieci dla SC to \emph{pierwsze maksimum lokalne
$J(a^*) > 0$} (atrakcyjne, sp\'ojne fazowo). Dla $g_0^e$:
\[
  a^*_0 = 7{,}75\ (\text{substr.}),
  \qquad J^*_0 = +7{,}5068\cdot 10^{-1}.
\]
Kolejne maksima: $a^*_1 = 14{,}00$, $a^*_2 = 20{,}50$ (struktura $a^*_{n+1}-a^*_n \approx 2\pi$).

\paragraph{$T_c$ z modelu XY na sieciach Bravais.}
Standardowe stosunki Monte Carlo $k_B T_c / J$ dla 3D XY
(aktualne warto\'sci literaturowe, Gottlob-Hasenbusch 1993, Janke~\& Nather 2005):

\begin{center}\footnotesize
\begin{tabular}{@{}lrrr@{}}
\toprule
Sie\'c     & $z$ & $k_B T_c/J$ & $T_c$ przy $J^* = 7{,}51\cdot10^{-1}$ \\
\midrule
SC (cubic)  & 6  & $2{,}2017$ & $+1{,}6528$ \\
BCC          & 8  & $2{,}9356$ & $+2{,}2037$ \\
FCC          & 12 & $4{,}4034$ & $+3{,}3055$ \\
2D square    & 4  & $0{,}8929$ (BKT) & $+0{,}6703$ \\
\bottomrule
\end{tabular}
\end{center}
($T_c$ w~jedn.~substratu; konwersja do kelvin\'ow: Appendix~\ref{app:P5-ps3}).

\label{app:P5-ps1-barrier}
\paragraph{Bariera vortex-ring: dow\'od topologicznej ochrony.}
Dla $d=3$ rozwicie globalnego windingu wymaga stworzenia
\emph{pier\'scienia wirowego} (vortex ring) o~promieniu $R$
(standardowy wynik London/Ginzburg-Landau).
Wz\'or~\eqref{eq:P5-Ering} daje:
\begin{center}\footnotesize
\begin{tabular}{@{}rrr@{}}
\toprule
$R$ (jedn.~substr.) & $E_\text{ring}$ & $E_\text{ring}/k_BT_c^\text{SC}$ \\
\midrule
$10$    & $+6{,}14\cdot10^{1}$ & $37{,}2$ \\
$100$   & $+1{,}70\cdot10^{3}$ & $1{,}03\cdot10^{3}$ \\
$1000$  & $+2{,}79\cdot10^{4}$ & $1{,}69\cdot10^{4}$ \\
\bottomrule
\end{tabular}
\end{center}
Bariera $E_\text{ring}/k_BT$ ro\'snie \emph{nieograniczenie} z~$R$ ---
dla makroskopowego $R\sim L_\text{sample}$ stan $n\neq 0$
jest \textbf{topologicznie chroniony}: $E_\text{barrier} \gg k_BT$
dla wszystkich $T<T_c$.

\label{app:P5-ps1-iff}
\paragraph{Werdykt IFF.}
Zerooporowo\'s\'c w~sieci soliton\'ow TGP \emph{istnieje wtw.}:
\begin{itemize}\itemsep0pt
  \item[(i)]   $J(a) > 0$ dla sta\l{}ej sieci $a$ (atrakcyjne sprz\k{e}\.zenie faz);
  \item[(ii)]  $T < T_c^{XY}(J,z,d)$ (poni\.zej progu koherencji);
  \item[(iii)] $\pi_1(U(1)) = \mathbb{Z}$ (netrywialna topologia param.~porz\k{a}dku).
\end{itemize}
\textbf{Konieczno\'s\'c:} zanik $J$ lub podgrzanie $T>T_c$ niszczy faz\k{e}.
\textbf{Wystarczalno\'s\'c:} (i)$+$(ii)$+$(iii) $\Rightarrow$
$\psi$ koherentna $+$ bariera vortex-ring $\Rightarrow$ pr\k{a}d persystentny
topologicznie chroniony (\eqref{eq:P5-Ering}).
\textbf{Status ps1: Hipoteza $\to$ Propozycja} (IFF numerycznie dla $g_0^e$;
analityczny dow\'od bariery vortex-ring jest klasyczny London/GL;
\emph{nowo\'sci\k{a} jest identyfikacja Friedel-oscylacji $J(a)$
z~samego ODE substratu TGP}).

% -------------------------------------------------------------------
\subsection{Problem P5~\#2: formu\l{}a skalująca $T_c(g_0, z, d)$}%
\label{app:P5-ps2}
% -------------------------------------------------------------------

\paragraph{Skan $g_0$ --- test uniwersalno\'sci $a^*$.}
Skrypt \texttt{ps2\_tc\_scaling.py} skanuje 18 warto\'sci
$g_0 \in [0{,}5,\,1{,}9]$, wyznaczaj\k{a}c dla ka\.zdej warto\'sci:
$A(g_0), \delta(g_0), a^*(g_0), J^*(g_0)$.
\emph{Kluczowy wynik} --- \textbf{$a^*$ jest uniwersalne}:
\begin{equation}
\label{eq:P5-astar}
  a^* \;=\; 7{,}725 \pm 0{,}016 \quad (\text{rozpr.~0{,}21\%, zakres } [7{,}711,\,7{,}766]).
\end{equation}
Uniwersalno\'s\'c p\l{}ynie \emph{z~linearyzacji ODE substratu}
wok\'o\l{} $g=1$: okres 2$\pi$ jest \textbf{niezale\.zny} od warto\'sci brzegowej $g_0$
(amplituda $A$ zale\.zy od $g_0$, ale okres nie).

\paragraph{Dopasowanie kwadratowe $J^* = C_0\cdot A^2$.}
Z~\eqref{eq:P5-Jpattern}: dwa ogony $\sim A\sin(r)/r$, sp\'ojne w~jednym
punkcie maksimum, daj\k{a} $J^* \propto A^2$. Numeryczny fit:
\begin{equation}
\label{eq:P5-C0}
  \boxed{\;\; J^* \;=\; C_0 \cdot A(g_0)^2, \qquad C_0 = 48{,}8222,
  \qquad \text{RMS~residuum} = 1{,}54\% \;\;}
\end{equation}
(wyznaczony bez wyrazu wolnego --- fit jednoparametrowy).
Residuum ma kszta\l{}t paraboliczny wokó\l{} $g_0 = 1$:
\[
  J^*/A^2 = 47{,}91 - 1{,}47(g_0-1) + 3{,}03(g_0-1)^2,
\]
wskazuj\k{a}c na ma\l{}e poprawki od~nielinearno\'sci ogona
(nie badamy analitycznie w~P5).

\paragraph{Kompletna formu\l{}a $T_c$.}
\begin{equation}
\label{eq:P5-Tc}
  \boxed{\;\;
   T_c(g_0, z, d)
   \;=\;
   k_d(z) \cdot C_0 \cdot A(g_0)^2 \cdot \frac{\Lambda_E}{k_B}
   \;\;}
\end{equation}
gdzie $k_d(z) = (k_BT_c/J)_\text{XY}(z,d)$ (patrz~tabelka ps1),
$A(g_0)$ z~numerycznego ODE, $C_0 = 48{,}8222$,
a~$\Lambda_E$ z~kalibracji (ps3).

\paragraph{Parametry zwi\k{e}kszaj\k{a}ce $T_c$.}
\begin{enumerate}[nosep]
  \item \textbf{Liczba koordynacyjna $z$}: $T_c \propto z$
    (liniowo, z~3D XY). FCC~(12) vs~SC~(6): czynnik 2.
  \item \textbf{Amplituda ogona $A(g_0)$}: $T_c \propto A^2$.
    Tau (A~=~0{,}96) vs~e~(A~=~0{,}125): czynnik $(0{,}96/0{,}125)^2 \approx 59$.
  \item \textbf{Wymiar $d$}: 3D~XY (MF) daje $T_c \approx 2{,}2J$;
    2D~BKT daje $T_c \approx 0{,}89 J$.
  \item \textbf{Sta\l{}a sieci $a \approx a^*$}: odchylenie od 7{,}725
    wprowadza kar\k{e} $J \to J\cos(\Delta a)$, geometrycznie.
\end{enumerate}

\paragraph{Status ps2: Szkic $\to$ Hipoteza.}
Formu\l{}a zamkni\k{e}ta w~4 parametrach ($g_0, z, d, \Lambda_E$).
Czeka na mikroskopowe wyprowadzenie $\Lambda_E$ z~ps3/P4 i~na weryfikacj\k{e}
w~ps4 na znanych materia\l{}ach.

% -------------------------------------------------------------------
\subsection{Problem P5~\#3: kalibracja substrat $\leftrightarrow$ SI}%
\label{app:P5-ps3}
% -------------------------------------------------------------------

\paragraph{Kalibracja d\l{}ugo\'sci: 1 jedn.~substratu $=$ promie\'n Bohra.}
Argument fizyczny: soliton $g_0(r)$ opisuje \emph{zlokalizowany stan
elektronowy} substratu TGP; naturaln\k{a} skal\k{a} elektronowej
lokalizacji w~atomie jest promie\'n Bohra $a_0 = 0{,}52918$~\AA.
St\k{a}d:
\begin{equation}
\label{eq:P5-length-calib}
  \boxed{\;\; 1\ \text{jedn.~substratu TGP} \;=\; a_0 \;=\; 0{,}52918\ \text{\AA} \;\;}
\end{equation}
Preferowane maksima TGP konwertuj\k{a} si\k{e} na:
\begin{center}\footnotesize
\begin{tabular}{@{}crrl@{}}
\toprule
$n$ & $a^*_n$ (substr.) & $a^*_n$ (\AA) & Obserwowane SC przy tej sta\l{}ej \\
\midrule
0 & $7{,}725$  & $4{,}088$  & \textbf{Al} (4{,}046~\AA, $T_c=1{,}175$~K) \\
1 & $14{,}008$ & $7{,}413$  & (poza prostymi metalami) \\
2 & $20{,}291$ & $10{,}738$ & (Ce/U zwi\k{a}zki, blisko) \\
3 & $26{,}575$ & $14{,}063$ & (superprzewodniki klatkowe) \\
\bottomrule
\end{tabular}
\end{center}

\paragraph{Idealny match: Al (4{,}046~\AA).}
Al ma $a/a_0 = 7{,}646$, residuum $|\Delta a^*| = 0{,}079$ (jedn.~substr.)
--- \emph{najmniejsze odchylenie w~ca\l{}ej bazie 74~materia\l{}\'ow}.
W~TGP nie~jest to przypadek: Al to najczystszy klasyczny BCS z~czystym
s-elektronem, wi\k{e}c sp\'o\l{}ka $g_0 = g_0^e$ dzia\l{}a idealnie.

\paragraph{Kalibracja energii: $\Lambda_E$ z~Al jako punkt referencyjny.}
Dla Al z~\eqref{eq:P5-Tc}:
\begin{align*}
  J^*(g_0^e) &= C_0 \cdot A_e^2 = 48{,}8222 \cdot (0{,}1246)^2 = 0{,}7578\ (\text{substr.}),\\
  T_c^\text{substr}(Al,SC) &= 2{,}20 \cdot J^* = 1{,}6685\ (\text{substr.}),\\
  \Lambda_E &= \frac{T_c^\text{obs}\cdot k_B}{T_c^\text{substr}}
             = \frac{1{,}175\,\text{K}\cdot k_B}{1{,}6685} = 60{,}687\ \mu\text{eV}.
\end{align*}
Dla kontekstu: pełna skala substratu $\sim E_\text{Hartree}\approx 27{,}2$~eV,
wi\k{e}c t\l{}umienie $E_\text{Hartree}/\Lambda_E \approx 4{,}5\cdot10^5$
\'---\ odpowiada to \emph{effective coupling} $1/\text{tlumienie}$
w~mikroskopowym opisie BCS.

\paragraph{Tabelka mapowania parametr\'ow efektywnych.}
\begin{center}\small
\begin{tabular}{@{}lll@{}}
\toprule
\textbf{Parametr materia\l{}u} & \textbf{$\leftrightarrow$ parametr TGP} & \textbf{Uwagi} \\
\midrule
Sta\l{}a sieci $a$~[\AA]        & $a/a_0$ w~jedn.~substratu; $a^*_n$ jak~\eqref{eq:P5-length-calib} & \textbf{MF} \\
$\omega_\text{Debye}$~[meV]    & skala pasma $\sim \Lambda_E$                & ok. 37~meV dla Al \\
$\lambda_\text{ep}$ (McMillan) & $A(g_0)^2 \cdot C_0 / \Lambda_E^\text{norm}$ & Al: 0{,}38; Pb: 1{,}55 \\
Koordynacja $z$                 & typ sieci Bravais (SC/BCC/FCC)              & $k_d(z)$ z~MC \\
$\xi_\text{Pippard}$             & $a_0 / A(g_0)$ (naiwnie)                    & \emph{krótki; faza sol.} \\
Gap $\Delta$                    & $E_\text{barrier}$ vortex-ring               & \eqref{eq:P5-Ering} \\
$\lambda_L$ (London)             & szeroko\'s\'c ekranowania substr.~U(1)      & dodatek O \\
\bottomrule
\end{tabular}
\end{center}

\paragraph{Test na 19~znanych SC.}
Przy za\l{}o\.zeniu uproszczonym $g_0 = g_0^e$ dla wszystkich materia\l{}\'ow
(ten sam substrat-elektron) oraz kary za niedopasowanie sta\l{}ej sieci
$\text{suppress} = \exp(-|\Delta a^*|^2/\sigma^2)$ ($\sigma = \pi/2$):
\[
  T_c^\text{TGP}(m) = 2{,}20 \cdot (z_m/6)\cdot C_0\cdot A_e^2
                    \cdot \text{suppress}(a_m)\cdot \Lambda_E/k_B.
\]
Korelacja log-log (Pearson) $r(\log T_c^\text{obs}, \log T_c^\text{pred}) = 0{,}306$.
Zgodno\'s\'c ilo\'sciowa ograniczona \emph{uproszczeniem jednego $g_0$}
--- rzeczywiste materia\l{}y maj\k{a} r\'o\.zne orbitale (d, f) z~innymi
warto\'sciami $g_0$, st\k{a}d rozrzut.

\paragraph{Status ps3: Program $\to$ Propozycja.}
\begin{itemize}[nosep]
  \item D\l{}ugo\'s\'c: \textbf{pe\l{}na} ($1$~substr.~$= a_0$, niezale\.zne od fitowania)
  \item Energia: \textbf{fenomenologiczna} ($\Lambda_E$ z~Al, jeden parametr)
  \item Zakres przewidywalnych $T_c$: od mK do dziesi\k{a}tek K (rz\k{a}d wielko\'sci)
\end{itemize}

% -------------------------------------------------------------------
\subsection{Problem P5~\#4: predykcja klas kandydat\'ow (TGP-score)}%
\label{app:P5-ps4}
% -------------------------------------------------------------------

\paragraph{Definicja TGP-score.}
Bez dopasowywania do bazy SC definiujemy \emph{a priori} rankujący
\begin{equation}
\label{eq:P5-score}
  \text{TGP-score}(m)
  \;=\;
  R_a(m)\cdot R_z(m)\cdot R_\text{orb}(m),
\end{equation}
gdzie:
\begin{itemize}\itemsep0pt
  \item $R_a(m) = \exp\!\bigl(-|a_m/a_0 - a^*_n|^2/\sigma_a^2\bigr)$, $\sigma_a = 1{,}5$,
    z~$n$ dobranym tak, aby minimalizowa\'c odchylenie;
  \item $R_z(m) = z_m/6$ (bonus koordynacyjny);
  \item $R_\text{orb}(m) = (A_\text{orb}/A_s)^2$ (amplituda ogona
    solitonu dla danej pow\l{}oki orbitalnej: $s, p, d, f$).
\end{itemize}
Wk\l{}ad $R_\text{orb}$ wzmacnia materia\l{}y z~d/f-orbitalami
(wi\k{e}ksze $A \Rightarrow T_c \propto A^2$ wi\k{e}ksze).

\paragraph{Wyniki na bazie 74~materia\l{}\'ow.}
Top-20 rankingu (\texttt{ps4\_candidate\_classes.py}):
\begin{center}\scriptsize
\begin{tabular}{@{}rlrrrr l@{}}
\toprule
Rank & Materia\l{} & $T_c$~[K] & $a$~[\AA] & orb & $z$ & TGP-score \\
\midrule
 1 & \textbf{Yb}      & $0{,}0$   & $3{,}883$ & f  & 12 & $110{,}18$ \\
 2 & La               & $6{,}0$   & $3{,}770$ & df & 12 & $100{,}32$ \\
 3 & \textbf{Ce}      & $0{,}0$   & $3{,}650$ & f  & 12 & $86{,}87$ \\
 4 & URu$_2$Si$_2$   & $1{,}4$   & $4{,}133$ & f  &  8 & $78{,}26$ \\
 5 & LaH$_{10}$      & $250{,}0$ & $3{,}320$ & df & 20 & $76{,}99$ \\
 6 & UBe$_{13}$      & $0{,}95$  & $10{,}260$& f  &  8 & $54{,}65$ \\
 7 & CeCoIn$_5$       & $2{,}3$   & $4{,}603$ & f  &  8 & $51{,}53$ \\
 8 & CaH$_6$         & $215{,}0$ & $3{,}650$ & d  & 20 & $35{,}32$ \\
 9 & YH$_6$           & $220{,}0$ & $3{,}580$ & d  & 20 & $31{,}80$ \\
10 & Ir                & $0{,}14$  & $3{,}839$ & d  & 12 & $26{,}04$ \\
11 & Th                & $1{,}38$  & $5{,}084$ & f  & 12 & $24{,}38$ \\
12 & \textbf{Y}       & $0{,}0$   & $3{,}647$ & d  & 12 & $21{,}10$ \\
13 & LaOFeAs           & $26{,}0$  & $4{,}035$ & d  &  8 & $19{,}07$ \\
14 & BaFe$_2$As$_2$    & $38{,}0$  & $3{,}963$ & d  &  8 & $18{,}69$ \\
15 & V$_3$Si           & $17{,}1$  & $4{,}725$ & d  & 14 & $17{,}60$ \\
16 & \textbf{FeTe}    & $0{,}0$   & $3{,}813$ & d  &  8 & $16{,}99$ \\
17 & LiFeAs            & $18{,}0$  & $3{,}775$ & d  &  8 & $16{,}40$ \\
18 & FeSe              & $8{,}0$   & $3{,}765$ & d  &  8 & $16{,}23$ \\
19 & V$_3$Ga           & $16{,}5$  & $4{,}817$ & d  & 14 & $14{,}42$ \\
20 & HgBa$_2$         & $135{,}0$ & $3{,}892$ & d  &  6 & $13{,}52$ \\
\bottomrule
\end{tabular}
\end{center}
\textbf{Walidacja:} 16~z~20 to znane SC; 4~pogrubione (Yb, Ce, Y, FeTe) to
kandydaci non-SC w~standardowych warunkach.

\paragraph{Konkrete predykcje eksperymentalne.}
\begin{itemize}\itemsep0pt
  \item \textbf{Yb, Ce, Y} (najwy\.zsze score non-SC, $f$ lub $d$):
    kandydaci na SC \emph{pod ci\'snieniem} --- Y znany jako SC pod ci\'sn.,
    Ce/Yb przewidywane z~modelu.
  \item \textbf{Clathrate / klatkowe (z=20)}: LaH$_{10}$, CaH$_6$, YH$_6$ ---
    wysokie $R_z = 3{,}33$ naturalnie wyja\'snia superhydrydy. Predykcja:
    inne klatkowe z~d/f-orbitalami b\k{e}d\k{a} SC pod ci\'snieniem.
  \item \textbf{Optymalne $T_c$ dla $a\approx 4{,}0-4{,}1$~\AA, d-orbitale,
    FCC/BCC}: czynnik $(A_d/A_s)^2 \approx 15$ vs~Al, oczekiwana
    $T_c \sim 15-30$~K.
  \item \textbf{Anti-predykcja: materia\l{}y z $a/a_0 \approx n\pi$ (zera $J(a)$)}
    nie b\k{e}d\k{a} SC. Mo (3{,}15~\AA) i~Tc (2{,}74~\AA) bliskie tego
    re\.zimu --- niskie $T_c$ mimo d-orbitali.
\end{itemize}

\paragraph{Klasy materia\l{}owe z~perspektywy TGP.}
Identyfikujemy \textbf{9 klas} zgodnych z~warto\'sci\k{a} TGP-score:
\begin{center}\scriptsize
\begin{tabular}{@{}clrrl@{}}
\toprule
Klasa & Opis & $N$ & $\langle T_c\rangle$ [K] & $\langle\text{score}\rangle$ \\
\midrule
I    & Ambient $s/p$ BCS              & 10 & $6{,}4$     & $4{,}0$  \\
II   & Transition-metal $d$           & 18 & $5{,}6$     & $7{,}1$  \\
III  & A15 intermetalic\'owe         &  6 & $15{,}7$   & $9{,}7$  \\
IV   & Cuprates $d$-orbital          &  4 & $115{,}5$  & $13{,}1$ \\
V    & Iron pnictide $d$              &  4 & $22{,}5$   & $17{,}6$ \\
VI   & Nickelate $d$                   &  2 & $27{,}5$   & $13{,}2$ \\
VII  & Super-hydryd (ci\'snienie)    & 10 & $100{,}6$  & $16{,}3$ \\
VIII & $f$-heavy fermion              &  7 & $1{,}9$     & $45{,}1$ \\
IX   & Kandydaci TGP (non-SC)         & 13 & $0{,}0$     & $19{,}0$ \\
\bottomrule
\end{tabular}
\end{center}
(Dane: \texttt{ps4\_results.txt}, sekcja~D.)

\paragraph{Status ps4: Program $\to$ Predykcja.}
Lista konkretnych kandydat\'ow do sprawdzenia eksperymentalnego
(non-SC, wysokie TGP-score):
\begin{center}
\begin{tabular}{@{}llrrr@{}}
\toprule
Materia\l{} & Sie\'c & $a$~[\AA] & orb & score \\
\midrule
Yb    & FCC         & $3{,}883$ & f  & $110{,}2$ \\
Ce    & FCC         & $3{,}650$ & f  & $86{,}9$  \\
Y     & HCP         & $3{,}647$ & d  & $21{,}1$  \\
FeTe  & tetragonal.  & $3{,}813$ & d  & $17{,}0$  \\
CaB$_6$ & cubic     & $4{,}149$ & sp & $7{,}7$   \\
\bottomrule
\end{tabular}
\end{center}

% -------------------------------------------------------------------
\subsection{Mapowanie BCS/Ginzburg-Landau~$\leftrightarrow$~substrat TGP}%
\label{app:P5-mapping}
% -------------------------------------------------------------------

\begin{center}\small
\begin{tabular}{@{}ll@{}}
\toprule
\textbf{BCS / Ginzburg-Landau} & \textbf{Substrat TGP} \\
\midrule
Para Coopera                        & Soliton $g_0$ z~dodatkow\k{a} faz\k{a} U(1) \\
Order parameter $\psi$              & $\psi = \sum_i g_0(r-R_i) e^{i\theta_i}$ \\
Gap $\Delta$                        & Bariera topologiczna rozwicia (vortex-ring, eq.~\eqref{eq:P5-Ering}) \\
$T_c$                                & Skala sprz\k{e}\.zenia $J(a^*) = C_0 A^2$, formu\l{}a~\eqref{eq:P5-Tc} \\
Coherence length $\xi$              & D\l{}ugo\'s\'c oscyluj\k{a}cego ogona $A/r$ (linearyz.~ODE) \\
Penetration depth $\lambda_L$       & Szeroko\'s\'c ekranu em w~substr.~U(1) (dod.~O) \\
Vortex / fluxoid                     & Soliton z~windingiem $n \in \pi_1(U(1)) = \mathbb{Z}$ \\
$\omega_\text{Debye}$                & Skala pasma $\sim \Lambda_E$ (fenomenologicznie) \\
$\lambda_\text{ep}$ (McMillan)      & $A(g_0)^2 \cdot C_0 / \Lambda_E^\text{norm}$ \\
Liczba koordynacyjna (BCS: $N(E_F)$) & Liczba koord.~sieci Bravais $z$ ($k_d(z)$ z~MC) \\
\bottomrule
\end{tabular}
\end{center}

\textbf{Kluczowa r\'o\.znica konceptualna:} BCS zak\l{}ada
ekspo-fall-off oddzia\l{}ywania poza $\xi$; TGP wyprowadza
\emph{Friedel-like oscylacje} $J(a) \sim \sin(a+\delta')/a^2$
z~samego ODE substratu --- nie z~fenomenologicznej funkcji
pomocniczej. To jest \emph{nowy j\k{e}zyk} opisowy, nie nowa fizyka.

% -------------------------------------------------------------------
\subsection{Tabela status\'ow twierdze\'n po~P5}%
\label{app:P5-status-table}
% -------------------------------------------------------------------

Konsolidacja status\'ow \textbf{twierdze\'n sektora nadprzewodnictwa TGP}
z~uwzgl\k{e}dnieniem wynik\'ow programu~P5. Kolumna ,,przed~P5''
odzwierciedla stan po~P4 (2026-04-19); kolumna ,,po~P5'' uwzgl\k{e}dnia
wyniki z~niniejszego dodatku.

\smallskip
\begin{center}\small
\setlength{\tabcolsep}{3pt}
\begin{longtable}{@{}p{5.2cm}lll p{3.2cm}@{}}
\toprule
\textbf{Twierdzenie / relacja} & \textbf{Przed~P5} & \textbf{Po~P5}
  & \textbf{Zmiana} & \textbf{Skrypt} \\
\midrule
\endfirsthead
\toprule
\textbf{Twierdzenie / relacja} & \textbf{Przed~P5} & \textbf{Po~P5}
  & \textbf{Zmiana} & \textbf{Skrypt} \\
\midrule
\endhead
\midrule
\multicolumn{5}{r}{\footnotesize\itshape ci\k{a}g dalszy\ldots} \\
\bottomrule
\endfoot
\bottomrule
\endlastfoot
%
\multicolumn{5}{l}{\textbf{A. IFF-warunek zerooporowo\'sci}}\\
\midrule
Sie\'c Bravais soliton\'ow jako prymityw
  & ---
  & \statuslabel{Def.}
  & \textbf{nowe}
  & ps1 \\[3pt]
Friedel-like oscylacje $J(a)$ z~ODE substr. (okres 2$\pi$)
  & ---
  & \statuslabel{Tw. [AN+NUM]}
  & \textbf{nowe}
  & ps1 \\[3pt]
IFF $(J{>}0)\land(T{<}T_c^{XY})\land(\pi_1{=}\mathbb{Z})$
  & ---
  & \statuslabel{Prop. [AN+NUM]}
  & \textbf{nowe}
  & ps1 \\[3pt]
Bariera vortex-ring $E_\text{ring}\propto J R\ln(R/\xi)$
  & \statuslabel{Tw. [AN]} (London)
  & \statuslabel{Tw. [AN]} (przeniesione)
  & brak (klasyczne)
  & ps1, dod.~T4 \\[3pt]
Zero $J(a_0) \Rightarrow a_0 = n\pi$
  & ---
  & \statuslabel{Tw. [NUM]}
  & \textbf{nowe} (4 cyfry)
  & ps1 \\[3pt]
%
\midrule
\multicolumn{5}{l}{\textbf{B. Formu\l{}a skaluj\k{a}ca $T_c$}}\\
\midrule
$J^* = C_0 A(g_0)^2$, $C_0 = 48{,}82$
  & ---
  & \statuslabel{Tw. [NUM]}
  & \textbf{nowe} (RMS 1{,}54\%)
  & ps2 \\[3pt]
$a^* = 7{,}725 \pm 0{,}016$ uniwersalne
  & ---
  & \statuslabel{Tw. [NUM]}
  & \textbf{nowe} (0{,}21\%)
  & ps2 \\[3pt]
$T_c = k_d(z)\cdot C_0 \cdot A^2 \cdot \Lambda_E/k_B$
  & ---
  & \statuslabel{Hip. [AN+NUM]}
  & \textbf{nowe} (formu\l{}a)
  & ps2 \\[3pt]
Uniwersalno\'s\'c $a^*$ z~linearyz.~ODE
  & ---
  & \statuslabel{Prop. [AN]}
  & \textbf{nowe}
  & ps2 \\[3pt]
%
\midrule
\multicolumn{5}{l}{\textbf{C. Kalibracja substrat~$\leftrightarrow$~SI}}\\
\midrule
1 jedn.~substratu $=$ $a_0$ (Bohr)
  & ---
  & \statuslabel{Def.}
  & \textbf{nowe} (kanoniczne)
  & ps3 \\[3pt]
$a^*_\text{TGP} = 4{,}088$~\AA
  & ---
  & \statuslabel{Prop. [NUM]}
  & \textbf{nowe} (Al match)
  & ps3 \\[3pt]
$\Lambda_E \approx 60{,}7$~$\mu$eV (z~Al)
  & ---
  & \statuslabel{Prop. [NUM]}
  & \textbf{nowe} (fenom.)
  & ps3 \\[3pt]
$\Lambda_E$ z~mikroskopowej teorii
  & ---
  & \statuslabel{Program}
  & otwarte
  & --- \\[3pt]
Test na 19 znanych SC (log-log $r = 0{,}306$)
  & ---
  & \statuslabel{Prop. [NUM]}
  & ilo\'sciowo mia\l{}kie
  & ps3 \\[3pt]
%
\midrule
\multicolumn{5}{l}{\textbf{D. Predykcja kandydat\'ow materia\l{}owych}}\\
\midrule
TGP-score = $R_a R_z R_\text{orb}$
  & ---
  & \statuslabel{Def.}
  & \textbf{nowe}
  & ps4 \\[3pt]
Top-20: 16 znanych SC + 4 kandydat\'ow
  & ---
  & \statuslabel{Tw. [NUM]}
  & \textbf{walidacja}
  & ps4 \\[3pt]
9 klas materia\l{}owych (I--IX)
  & ---
  & \statuslabel{Prop.}
  & \textbf{nowe}
  & ps4 \\[3pt]
Yb, Ce, Y kandydaci pod ci\'snieniem
  & ---
  & \statuslabel{Pred.}
  & \textbf{nowe}
  & ps4 \\[3pt]
Anti-predykcja: $a/a_0 \approx n\pi \Rightarrow$ non-SC
  & ---
  & \statuslabel{Pred.}
  & \textbf{nowe}
  & ps4 \\[3pt]
Clathrate z~d/f $+$ P $\Rightarrow$ super-hydryd
  & ---
  & \statuslabel{Pred.}
  & \textbf{nowe}
  & ps4 \\[3pt]
%
\midrule
\multicolumn{5}{l}{\textbf{E. Bilans sektora nadprzewodnictwa}}\\
\midrule
Liczba wolnych parametr\'ow teorii
  & --- (brak teorii P5)
  & \textbf{2} ($g_0^e$, $\Lambda_E$)
  & brak (minimum)
  & --- \\[3pt]
Liczba predykcji ilo\'sciowych $T_c$
  & 0
  & \textbf{19} (baza testowa)
  & \textbf{+19}
  & ps3 \\[3pt]
Nazwane sta\l{}e TGP (z~P5)
  & brak
  & \textbf{2} ($C_0{=}48{,}82$, $a^*{=}7{,}725$)
  & \textbf{+2}
  & ps1, ps2 \\
%
\end{longtable}
\end{center}

\paragraph{Podsumowanie zmian (5-wierszowe).}
\begin{enumerate}[nosep]
  \item Friedel-like oscylacje $J(a)$ \emph{z~ODE substratu}:
    \textbf{nowe twierdzenie [NUM]} (4 cyfry, okres $2\pi$).
  \item IFF-warunek zerooporowo\'sci: \textbf{Propozycja [AN+NUM]}
    (wariant A topologiczny, baza $\beta$-IFF).
  \item Formu\l{}a $T_c = k_d(z) C_0 A^2 \Lambda_E/k_B$:
    \textbf{Hipoteza [AN+NUM]} (RMS 1{,}54\%).
  \item Kalibracja: 1 jedn.~substr. $=$ $a_0$ $\Rightarrow$ $a^*_\text{TGP} = 4{,}088$~\AA{}
    (Al match, $T_c = 1{,}18$~K).
  \item TGP-score ranking: 16/20 znane SC, kandydaci non-SC
    (Yb, Ce, Y, FeTe, CaB$_6$) do sprawdzenia pod ci\'snieniem.
\end{enumerate}

% -------------------------------------------------------------------
\subsection{Otwarte pytania po~P5}%
\label{app:P5-open}
% -------------------------------------------------------------------

\begin{enumerate}[nosep]
  \item \textbf{Mikroskopowe $\Lambda_E$.} Obecnie $\Lambda_E \approx 60{,}7$~$\mu$eV
    z~Al (fenomenologicznie). Czy~$\Lambda_E$ wynika z~innych sta\l{}ych
    TGP (np. $\alpha_3$, $\rho_0^*$, skale ERG)? --- Program P5b.
  \item \textbf{Mapa $g_0 \leftrightarrow$ orbitale pow\l{}okowe.}
    Dla~rzeczywistych materia\l{}\'ow $g_0$ zale\.zy od symetrii orbitalu
    (s/p/d/f). Obecnie u\.zywamy $g_0^e$ (s-like); potrzebna~rozszerzona
    baza soliton\'ow substratowych dla d/f (moduł P6 kandydat).
  \item \textbf{Korekta drugiego rz\k{e}du do $C_0$.} Fit kwadratowy
    $J^*/A^2 = 47{,}91 - 1{,}47(g-1) + 3{,}03(g-1)^2$ sugeruje
    systematyczne odchylenie od jednorodnego $C_0 = 48{,}82$.
    Wsp\'o\l{}czynnik $3{,}03$ pr\'obujemy dopasowa\'c do kombinacji
    $\pi^2/3 - \alpha_3$ lub~podobnych.
  \item \textbf{Sieci niesferyczne (A15, perovskite).} Formu\l{}a \eqref{eq:P5-Tc}
    zak\l{}ada czysto-kuboidalne sieci Bravais. A15~(Nb$_3$Sn, V$_3$Si) maj\k{a}
    kana\l{}y $1D$; cuprates maj\k{a} p\l{}aszczyzny CuO$_2$ ($d=2$).
    Rozszerzenie TGP na niesferyczne geometrie jest osobnym problemem.
  \item \textbf{Dow\'od $\beta$-IFF.} Obecnie IFF pokazana \emph{numerycznie}
    + wariant A topologiczny dla wystarczalno\'sci. Pe\l{}ny dow\'od
    konieczno\'sci (,,je\'sli nie~wszystkie (i)(ii)(iii), to~SC nie~istnieje'')
    wymaga eliminacji scenariuszy neklasycznych (fluktuacje kwantowe
    przy $T \to 0$).
  \item \textbf{Eksperymentalna weryfikacja kandydat\'ow.}
    \textbf{Yb, Ce, Y} pod~ci\'snieniem $P > 10$~GPa: czy~obserwuje si\k{e}
    $T_c > 1$~K? Yb nigdy nie by\l{} zweryfikowany pod~bardzo wysokim P
    poza zakresem magnetycznym.
\end{enumerate}

\medskip
\textbf{Skrypty i~dane (P5):}
\begin{itemize}[nosep]
\item \texttt{research/superconductivity\_closure/ps1\_zero\_resistance\_condition.py} --- IFF-warunek, Friedel $J(a)$, $T_c$ XY
\item \texttt{research/superconductivity\_closure/ps2\_tc\_scaling.py} --- skan $g_0$, $J^* = C_0 A^2$, fit $C_0 = 48{,}82$
\item \texttt{research/superconductivity\_closure/ps3\_effective\_params\_map.py} --- kalibracja $a_0$ + $\Lambda_E$, test 19 SC
\item \texttt{research/superconductivity\_closure/ps4\_candidate\_classes.py} --- TGP-score, baza 74 mat., 9 klas
\item \texttt{research/superconductivity\_closure/README.md} --- skonsolidowane wyniki P5
\item \texttt{research/superconductivity\_closure/ps\{1,2,3,4\}\_results.txt} --- pe\l{}ne output'y
\end{itemize}

\paragraph{Powiązania z rdzeniem TGP.}
\begin{itemize}[nosep]
  \item sek.~\ref{sec:pole} / dodatek~B --- ODE substratu ($\alpha{=}1, d{=}3$): baza solitonowa.
  \item dodatek~J (\ref{app:J-profil}, \ref{app:J-skalowanie}) --- ogon masy, amplituda $A_\text{tail}$ dla~soliton\'ow.
  \item dodatek~T4 (\ref{app:T4-bounce}) --- bounce topology, $\pi_1(U(1)) = \mathbb{Z}$.
  \item dodatek~O (\ref{app:u1_formal}) --- U(1)-formalizacja, faza $\psi = |\psi|e^{i\theta}$.
  \item dodatek~R2 (\ref{app:R2-qk-z3}) --- dynamical Z$_3$, analogia dla 3-band SC / $d$-wave.
  \item Program~P4 (\ref{app:P4}) --- baza soliton\'ow $g_0^{e,\mu,\tau}$.
\end{itemize}
